\documentclass[11pt]{article}
\usepackage[T1]{fontenc}
\usepackage{lmodern,microtype,amsmath,amsthm,amssymb,booktabs}
\usepackage[margin=1in]{geometry}
\usepackage[colorlinks=true,urlcolor=blue,linkcolor=black,citecolor=black]{hyperref}
\hypersetup{pdftitle={The word aabbaba forces periodicity}}
\newtheorem{theorem}{Theorem}
\newtheorem{lemma}[theorem]{Lemma}
\newtheorem{corollary}[theorem]{Corollary}
\newcommand{\eps}{\varepsilon}
\title{The word \(\mathtt{aabbaba}\) forces periodicity}
\author{Aditya Balaji}
\date{Research manuscript, 24 November 2026}
\theoremstyle{remark}\newtheorem*{remark}{Remark}
\begin{document}
\maketitle
\begin{abstract}
We prove that \(XXYYXYX=UUVVUVU\) in a free monoid implies either \(X=U,\ Y=V\), or all four words commute. No length ordering or nonerasing hypothesis is required. Two overlapping factorizations force an intervening word to be a cube. Transporting this cube and then a fifth power through the equation produces a common commuting word. We also determine exactly which letter positions are forced to agree at every choice of lengths. A Lean formalization checks the unrestricted theorem. We also prove that \(ab^3a\) forces periodicity under Holub's length condition, completing all sixteen tuples with exponents at most one (also checked in Lean), and we give explicit nonforcing families for all exponents together with a conjectured complete classification that matches an exact search over all 625 tuples with exponents at most four.
\end{abstract}
\section{Statement}
Holub's problem at WORDS 2025 asks which patterns \(a^i b a^j b a^k b a^\ell\) force periodicity when \(V\) is at least as long as each of \(X,Y,U\) \cite[pp.~3--5]{words}. The simplest unknown case identified there has \((i,j,k,\ell)=(2,0,1,1)\). The slides use finite words and a morphism formulation without an explicit nonerasing requirement. Our theorem covers both conventions.
\begin{theorem}\label{main}
Over any alphabet, all finite words satisfying \(XXYYXYX=UUVVUVU\) commute pairwise or satisfy \(X=U,\ Y=V\).
\end{theorem}
\begin{lemma}[Transport of a power]\label{transport}
If \(P^m Z=Z Q\) with \(m\ge1\), there is a word \(K\) with \(PZ=ZK\) and \(Q=K^m\). In particular \(|K|=|P|\).
\end{lemma}
\begin{proof}
If \(P=\eps\), length gives \(Q=\eps\). Otherwise induct on \(|Z|\). At length zero take \(K=P\). If \(Z=cZ'\), the first letter of \(P\) is \(c\); write \(P=cR\). Cancelling \(c\) in the hypothesis gives \((Rc)^m Z'=Z'Q\), since \((cR)^m c=c(Rc)^m\). Induction gives \(RcZ'=Z'K\) and \(Q=K^m\). Prepending \(c\) gives \(PZ=ZK\).
\end{proof}
If \(P,Q\) commute with a nonempty \(H\), then \(PQ=QP\). Indeed both concatenations commute with every \(H^n\); for \(|H^n|\ge |PQ|\) both are prefixes of \(H^n\) of the same length.

\section{Proof of Theorem~\ref{main}}
Exchange the two substitutions if necessary to arrange \(|Y|\le |V|\). Length gives \(4(|X|-|U|)=3(|V|-|Y|)\), hence
\[
 |X|-|U|=3t,\qquad |V|-|Y|=4t,\qquad t\in\mathbb Z_{\ge0}.
\]
If \(t=0\), corresponding boundaries coincide and the corresponding words agree. Henceforth \(t>0\).

At each boundary compare the two prefixes of the common expanded word. The diagram shows the left prefix length minus the right, in units of \(t\).
\[
\begin{array}{c|rrrrrrrr}
\text{boundary}&0&1&2&3&4&5&6&7\\ \hline
\text{left block just read}&-&X&X&Y&Y&X&Y&X\\
\text{right block just read}&-&U&U&V&V&U&V&U\\
\text{signed difference}/t&0&3&6&2&-2&1&-3&0\\
\text{intervening word}&-&A&B&C&D&E&F&-
\end{array}
\]
At boundary three the left prefix is the right prefix followed by \(C\); at boundary four the right prefix is the left prefix followed by \(D\). Appending the next blocks and cancelling common prefixes gives
\[
 X=UA,\quad AX=UB,\quad BY=VC,\quad V=CYD,\quad
 X=DUE,\quad V=EYF,\quad X=FU .
\]
Only the order of the corresponding boundaries matters; an intervening word may cross other block boundaries.

\paragraph{The forced cube.}
Compare \(CYD=EYF\). Since \(|C|=2t,\ |E|=t\), write \(C=EG\). Since \(|F|=3t,\ |D|=2t\), comparison at the right gives \(F=HD\) with \(|H|=t\). Now \(DUE=X=HDU\). Its first \(2t\) letters form \(D\). Its first \(t\) letters are \(H\), and its next \(t\) letters are the first \(t\) of \(D\), again \(H\). Thus
\[
 D=H^2,\qquad F=H^3,\qquad HU=UE.
\]
In particular \(X=H^3U=UE^3\). Since \(X=UA\), cancellation gives \(A=E^3\).

\paragraph{The first transport.}
The equation \(AX=UB\) becomes \(AUA=UB\). Since \(|B|=6t\), it ends in \(A\); write \(B=B_0 A\). Right cancellation gives \(E^3 U=U B_0\). Lemma~\ref{transport} gives \(EU=UK\) and \(B_0=K^3\), with \(|K|=t\). But \(BY=VC\), and \(V=EYF\). Comparing their first \(t\) letters gives \(K=E\). Therefore \(EU=UE=HU\); right cancellation gives \(H=E\). Consequently \(HU=UH\) and \(B=H^6\).

\paragraph{The second transport.}
Substitute \(C=HG,\ V=HYH^3\) into \(BY=VC\) and cancel the first \(H\):
\[
 H^5Y=YH^4G .
\]
Lemma~\ref{transport} gives \(HY=YL\) and \(L^5=H^4G\), with \(|L|=t\). Their first \(t\) letters give \(L=H\). Thus \(Y\) commutes with \(H\). So do \(U\), \(X=UH^3\), and \(V=HYH^3\). Since \(|H|=t>0\), the common-centralizer observation proves the claim. \qed

\section{All fixed-length solutions}
\begin{corollary}\label{gcd}
Fix nonnegative lengths \(x,y,u,v\) with \(4x+3y=4u+3v\). For \(x=u\) the solutions are precisely the identical substitutions. Otherwise set \(d=\gcd(x,y,u,v)>0\). The solutions are exactly
\[
 (X,Y,U,V)=(P^{x/d},P^{y/d},P^{u/d},P^{v/d}),\qquad |P|=d .
\]
The graph identifying corresponding letter positions has exactly \(d\) components, indexed by the local position modulo \(d\).
\end{corollary}
\begin{proof}
Commuting words are powers of a common primitive word: for a pair \(AB=BA\) with \(0<|A|\le |B|\), write \(B=AC\) and cancel to get \(AC=CA\), reducing total length. Uniqueness of primitive roots lets us apply this fact to each pair with a fixed nonempty word; see also \cite[Lemma 1]{hs}. The common root length divides \(d\); repeat the root to length \(d\) to obtain \(P\). Conversely both expansions are powers of \(P\) of equal length.

Assign each component its own symbol. This is a solution, so all positions of equal residue carry the same symbol. Distinct residues cannot be identified because each occurrence starts at an expanded position divisible by \(d\). Every residue occurs. For \(t>0\), the same \(d\) is \(\gcd(u,y,t)\), since \(\gcd(3,4)=1\).
\end{proof}

For any fixed pattern and lengths, build this graph by expanding into references rather than letters. All solutions are exactly the component labelings. A mismatch between symbolic \(AB\) and \(BA\) survives a binary letter map separating the two offending symbols. In the target case unequal corresponding lengths already guarantee distinct substitutions. In general, preserving both noncommutation and unequal substitutions needs at most two unequal-symbol constraints: any differing-length comparison survives automatically. A graph with at most two edges between distinct symbols is bipartite. This proves binary completeness at the same lengths. Scaling lengths by \(k\) replaces the graph with \(k\) independent copies, one per residue inside a scaled letter slot.

\section{Other exponent quadruples}
\begin{corollary}
For \(d\ge1\) and \(r,s\ge0\), the quadruple
\[
 (2d+r,\ s+r,\ d+s+r,\ d+s)
\]
forces periodicity without length restrictions, as does its reversal.
\end{corollary}
\begin{proof}
Replace \(a\) by \(a^d\) and \(b\) by \(a^r b a^s\) in the base pattern. Apply Theorem~\ref{main} to the four composite blocks \(X^d,X^rYX^s,U^d,U^rVU^s\). Equality of substitutions gives \(X=U\) by comparing equal-length prefixes of equal \(d\)-th powers, then \(Y=V\) by cancellation. Otherwise the four blocks are powers of a common primitive word, unless all are empty, when the conclusion is immediate. By uniqueness of primitive roots, \(X,U\) themselves are powers of this word (zeroth powers allowed). Removing their powers at the ends of the other blocks gives the same for \(Y,V\). Reversing each word proves the reversal assertion. The padding closure for \(d=1\) is \cite[Lemma 26]{hs}; the forcing seed here is Theorem~\ref{main}.
\end{proof}
For an infinite negative family satisfying the maximum-length condition, take \(a^m b^3\) with \(m\ge3\). Put \(T=0^{3m}\) and
\[
 X=(T1)^3,\quad Y=0^m,\quad U=000,\quad V=(1T)^m .
\]
The identity \((T1)^{3m}T=T(1T)^{3m}\) verifies the equation. Here \(X\ne U\), \(X,Y\) do not commute, and \(|V|=m(3m+1)\ge 3(3m+1)=|X|\). This specializes the construction in the proof of \cite[Theorem 30]{hs}.

\paragraph{Small exponents.}
For \(i,j,k,\ell\in\{0,1\}\), the table gives eleven nonforcing cases with nonempty binary certificates. Four other tuples force periodicity: \(0011,1100,0101,1010\). Their words are \(bbaba,ababb,babba,abbab\), covered up to reversal and letter exchange by \cite[Lemmas 14--15]{hs} for nonempty images (equal corresponding lengths give identical substitutions). If a substitution erases a letter, it is periodic. If just one substitution is periodic, \cite[Theorem 1]{periodic} requires the equality word to be unary or a power of a word having a letter occurring precisely once. Our four patterns are primitive by their coprime multiplicities 2 and 3, and each letter occurs at least twice. If both substitutions are periodic, their equal nonempty image has one primitive root; equality with the empty image forces all images empty. This covers erasures.

The tuple \(1001\), or \(ab^3a\), forces periodicity under the length condition by Theorem~\ref{t1001} below, although the odd-exponent example in \cite[conclusion]{hh} shows it fails without it. Thus all sixteen tuples in the cube \(\{0,1\}^4\) are classified.
\begin{center}\small\ttfamily
\begin{tabular}{lllll}\toprule
tuple & X & Y & U & V\\\midrule
0000&1&0&0&0\\
0001&0001&000&1&0000\\
0010&0010&000&1&0000\\
0100&0100&000&1&0000\\
0110&1010101&010&1&0101010\\
0111&10&0&0&01\\
1000&1000&000&1&0000\\
1011&101&0&1&010\\
1101&101&0&1&010\\
1110&10&0&1&00\\
1111&1010101&0&1&010101010\\\bottomrule
\end{tabular}
\end{center}
Without length restrictions, existence of a counterexample for any fixed pattern is decidable in principle by Makanin's theorem \cite{mak}: over a binary alphabet a disequality is a finite disjunction of proper-prefix and first-mismatching-letter cases, each expressible with word equations. Empty variables can be guessed separately. This gives no direct encoding here of the extra length comparisons. The fixed-length method and Corollary~\ref{gcd} do not classify all exponent quadruples.

% Section for the Holub manuscript: the tuple (1,0,0,1).
% Uses \eps for the empty word, as in the main manuscript.
\section{The tuple \((1,0,0,1)\)}
\begin{theorem}\label{t1001}
Over any alphabet, if \(XYYYX=UVVVU\) and \(|X|,|Y|\le|V|\), then \(X=U,\ Y=V\), or all four words commute pairwise. Hence \(ab^3a\) is periodicity forcing in Holub's sense.
\end{theorem}
Without \(|X|\le|V|\) the statement fails \cite{hh}; the proof uses that inequality exactly once, in Step~4. Write \(|Z|\) for length, \(Z[i{:}j]\) for the factor of positions \(i,\dots,j-1\) (from \(0\)), and \(\operatorname{rot}_r(w)=w[r{:}]\,w[{:}r]\). We use three standard facts.
\begin{itemize}
\item[(F1)] If \(ZW=WZ\) and \(W\ne\eps\), then \(Z\) is a power of the primitive root of \(W\).
\item[(F2)] If \(\operatorname{rot}_c(\tau)=\tau\) and \(\tau=\rho^k\) with \(\rho\) primitive, then \(|\rho|\) divides \(c\).
\item[(F3)] (Fine--Wilf) A word with periods \(p,q\) and length at least \(p+q-\gcd(p,q)\) has period \(\gcd(p,q)\).
\end{itemize}

\begin{lemma}\label{lemC}
If \(|B|=|D|=s>0\) and \(K\,BDB=DBD\,K\), then \(B=D\).
\end{lemma}
\begin{proof}
If \(|K|\ge3s\) then \(K=DBD\,K_1\) and \(K_1\) satisfies the same equation, so assume \(r=|K|<3s\). Then \(DBD=KL\), \(BDB=LK\), so \(BDB=\operatorname{rot}_r(DBD)\). If \(r\in\{0,s,2s\}\), compare blocks: \(BDB\) equals \(DBD\), \(BDD\) or \(DDB\), and in each case one block gives \(B=D\).

Otherwise \(r=qs+t\) with \(0<t<s\). For a word \(Z\) of length \(s\) put \(Z'=Z[{:}t]\), \(Z''=Z[t{:}]\). With \(\alpha=\alpha_1\alpha_2\alpha_3=\operatorname{rot}_{qs}(DBD)\), the blocks of \(BDB=\operatorname{rot}_t(\alpha)\) are \(\alpha_j''\alpha_{j+1}'\), indices mod 3.

\emph{Case \(q=0\), \(\alpha=DBD\):} \(B=D''B'\), \(D=B''D'\), \(B=D''D'\). The first and third give \(B'=D'\). Hence \(B=\operatorname{rot}_t(D)\) and \(D=\operatorname{rot}_t(B)\), so \(\operatorname{rot}_{2t}(B)=B\). Let \(\rho\) be the primitive root of \(B\), \(p=|\rho|\); by (F2) \(p\mid 2t\). If \(p\mid t\), then \(D=\operatorname{rot}_t(B)=B\). Otherwise \(p\) is even and \(t\equiv p/2\pmod p\), so \(t\ge p/2\). Now \(B'=D'\) reads \(B[i]=B[(i+t)\bmod s]\) for \(i<t\). For \(i<p/2\) this gives \(\rho[i]=\rho[i+p/2]\), so \(\rho\) is a square, contradicting primitivity.

\emph{Case \(q=1\), \(\alpha=BDD\):} \(B=B''D'\), \(D=D''D'\), \(B=D''B'\). Prefixes of length \(s-t\) of the first two show that \(B\) and \(D\) both have period \(t\). The suffixes of length \(t\) of the first and third give \(D'=B[s-t{:}]=B'\). Words of length \(s\) with period \(t<s\) are determined by their first \(t\) letters, so \(B=D\).

\emph{Case \(q=2\):} from \(DBD=KL\), \(BDB=LK\) we get \(L\,DBD=BDB\,L\) with \(|L|=3s-r\in(0,s)\). This is the case \(q=0\) with the roles of \(B\) and \(D\) exchanged, so again \(B=D\).
\end{proof}

\begin{proof}[Proof of Theorem~\ref{t1001}]
\emph{Step 1 (lengths).} Put \(x=|X|\) and so on. Then \(2x+3y=2u+3v\) and \(y\le v\) give \(x-u=3s\), \(v-y=2s\) with \(s\ge0\). If \(s=0\), corresponding block boundaries coincide, so \(X=U\) and \(Y=V\). Let \(s>0\). The hypothesis \(x\le v\) becomes
\begin{equation}\label{len}
 y\ \ge\ u+s .
\end{equation}

\emph{Step 2 (boundaries).} Comparing the two factorizations at their block boundaries (left prefix minus right prefix: \(3s,\ s,\ -s,\ -3s,\ 0\)) gives words with \(|A|=|E|=3s\), \(|B|=|D|=s\) and
\[
 X=UA=EU,\qquad AY=VB,\qquad V=BYD,\qquad YE=DV .
\]
Substituting \(V\) and cancelling (\(A\) begins with \(B\), \(E\) ends with \(D\)), we write \(A=BA'\), \(E=E'D\), \(P=DB\) with \(|A'|=|E'|=|P|=2s\), and get
\begin{equation}\label{core}
 A'Y=YP,\qquad YE'=PY,\qquad X=UBA'=E'DU .
\end{equation}

\emph{Step 3 (it suffices that \(YP=PY\)).} Suppose \(YP=PY\). Then \(A'Y=PY\) and \(YE'=YP\), so \(A'=E'=P\) and \eqref{core} becomes \(U\,BDB=DBD\,U\). Lemma~\ref{lemC} gives \(B=D\), hence \(P=B^2\) and \(UB^3=B^3U\). Let \(\rho\) be the primitive root of \(B\ne\eps\). By (F1), \(U\) and \(Y\) (which commutes with \(B^2\)) are powers of \(\rho\), and so are \(X=UB^3\) and \(V=BYB\). All four words commute.

\emph{Step 4 (\(YP=PY\)).} If \(y\ge2s\), then \(YE'=PY\) shows \(Y[{:}2s]=P\), and \(A'Y=YP\) shows \(A'=Y[{:}2s]=P\); hence \(PY=YP\).

Let \(y<2s\) and \(\ell=2s-y\). By \eqref{len}, \(s\le y\), \(u<s\) and \(0<\ell\le s\le y\). From \(YE'=PY\) the word \(Y\) is a prefix of \(P\): \(P=YR\) with \(|R|=\ell\), and \(E'=RY\). From \(A'Y=YYR\) we get \(A'=Y\,Y[{:}\ell]\) and \(Y=Y[\ell{:}]R\). So \(Y\) has period \(\ell\) and ends with \(R\). Also \(D=Y[{:}s]\) and \(B=Y[s{:}]R\). Hence \eqref{core} reads
\[
 \mathcal W:=R\,Y\,Y[{:}s]\,U\;=\;U\,Y[s{:}]\,R\,Y\,Y[{:}\ell].
\]
Let \(\delta=u+y-s=|U\,Y[s{:}]|\), so \(0\le\delta<y\). If \(\delta=0\), then \(y=s=\ell\) and \(R=Y\). Otherwise \(RY\) occurs in \(\mathcal W\) at positions \(0\) and \(\delta\), so \(\mathcal W[i]=\mathcal W[i+\delta]\) for \(i<2s\). In particular \(Y=\mathcal W[\ell{:}2s]\) has periods \(\delta\) and \(\ell\). Since \(y\ge u+s=\delta+\ell\) by \eqref{len}, (F3) gives \(Y\) period \(g=\gcd(\delta,\ell)\).

Next, \(\mathcal W[2s{:}3s]=Y[{:}s]\). For \(j<\min(s,\delta)\), the index \(2s+j-\delta\) lies in \([0,2s)\), so
\[
 Y[j]=\mathcal W[2s+j]=\mathcal W[2s+j-\delta]=Y[y-\delta+j].
\]
Both indices are below \(y\) and \(g\mid\delta\), so with \(c=y\bmod g\) and \(\tau=Y[{:}g]\) we get \(\tau=\operatorname{rot}_c(\tau)\), using \(j<g\le\min(s,\delta)\) and \(y\ge2g\). By (F2) the primitive root \(\rho\) of \(\tau\) has \(|\rho|\mid c\) and \(|\rho|\mid g\), hence \(|\rho|\) divides \(y\) and \(\ell\). Thus \(Y=\rho^{y/|\rho|}\), and its suffix \(R\) is \(\rho^{\ell/|\rho|}\). So \(YR=RY\), and \(P=YR\) commutes with \(Y\).

Steps 3 and 4 prove the theorem.
\end{proof}

\section{A conjectured complete classification}
For \(n,m\ge1\) put \(X_n=a^nba^n\), \(Y_n=(ba^{2n})^{n-1}b\), \(V_n=Y_na^nY_n\). Since \(X_n^n=a^nY_na^n\), substituting \((X,Y,U,V)=(X_n,Y_n,a,V_n)\) makes both sides of the equation equal for the tuples
\[
(n,0,n,n),\ (n,n,0,n),\ (0,n,2n,0),\ (0,2n,n,0),\ (0,2n,0,n),\ (n,0,2n,0);
\]
for instance both sides of \((n,0,n,n)\) expand to \(a^nY_na^nY_nY_na^nY_na^nY_na^nY_na^n\). Here \(|V_n|=2|Y_n|+n\ge 2n+1=|X_n|\) and \(X_n\) does not commute with \(a\). The tuples \((i,m,m,m-i)\) are nonforcing via \(X=a,\ Y=b,\ U=\eps,\ V=a^iba^{m-i}\), whose image is \(V^3\).

\begin{theorem}\label{nonforcing}
The tuple is not periodicity forcing if it lies in one of: \(A=\{(i,m,m,m-i)\}\); \(B=\{(m,0,0,0),(0,0,0,m)\}\); \(C=\{(n,0,n,n),(n,n,0,n)\}\); \(D=\{(0,n,2n,0),(0,2n,n,0),(0,2n,0,n),(n,0,2n,0)\}\); or \(\{0010,0100,0110,1111\}\).
\end{theorem}
\begin{proof}
Families \(C,D\) and \(A\) are above. Family \(B\) with \(m\ge3\) is the construction of Section~4 and its reversal; \(m=1,2\) and the four sporadic tuples have explicit binary witnesses in the companion files (\texttt{families.py}).
\end{proof}

\paragraph{Conjecture.} Every other tuple forces periodicity. At fixed lengths, union-find on letter positions yields the most general solution, of which every solution is a letter image; hence a counterexample exists at those lengths if and only if the generic solution is one. For all \(625\) tuples with exponents at most \(4\) and all lengths with \(|V|\le150\), the nonforcing tuples are exactly the \(43\) predicted by Theorem~\ref{nonforcing}; the largest witness needed has \(|V|=60\). Spot checks to \(|V|\le250\) agree for larger exponents. Theorems~\ref{main} and~\ref{t1001} prove the conjecture for \(2011\) and \(1001\).

\section{Reductions toward the full classification}
Write \(w=a^iba^jba^kba^\ell\), \(S=i+j+k+\ell\), and let \(g,h\) be the morphisms \(g(a)=X,\ g(b)=Y,\ h(a)=U,\ h(b)=V\), so the equation reads \(g(w)=h(w)\). Let \(\mathrm{Eq}(g,h)=\{z: g(z)=h(z)\}\). A binary morphism is \emph{periodic} if its two images commute (this includes erasing morphisms); a binary morphism is injective exactly when it is not periodic.

\begin{lemma}[Periodic morphisms]\label{perred}
If \(g\) or \(h\) is periodic, then \(X=U,\ Y=V\), or all four words commute, or the tuple lies in family \(A\), family \(B\), or \(\{0100,0010\}\).
\end{lemma}
\begin{proof}
If both are periodic and \(g(w)=h(w)\ne\eps\), the common word has a unique primitive root, which is then the root of all nonempty images; so all four words commute. If \(g(w)=h(w)=\eps\), all images of letters occurring in \(w\) are empty. If exactly one is periodic, \cite[Theorem 1]{periodic} gives, up to exchanging \(a\) and \(b\), \(\mathrm{Eq}(g,h)=a^*\) or \((a^rba^s)^*\). As \(w\) contains \(b\), either \(w=(a^rba^s)^3\), giving \((r,r+s,r+s,s)\in A\), or \(w=(b^rab^s)^c\) with \(c(r+s)=3\): for \(c=1\) the tuple has a single \(a\), so it is \(1000,0001\in B\) or \(0100,0010\); for \(c=3\), \(w\in\{(ab)^3,(ba)^3\}\subseteq A\). The tuple \(0000\) also lies in \(A\).
\end{proof}

\begin{theorem}[Balanced prefixes]\label{balanced}
Let \(S=3m\) with \(m\ge1\), and suppose \(i\le m\le i+j\) or \(i+j\le 2m\le i+j+k\). If the tuple is not in \(A\cup C\cup D\), it forces periodicity, even without the length condition.
\end{theorem}
\begin{proof}
By Lemma~\ref{perred} assume \(g,h\) injective, and \(g\ne h\) (otherwise \(X=U,\ Y=V\)). Lengths give \(S(x-u)=3(v-y)\), so for every prefix \(p\) of \(w\),
\[|g(p)|-|h(p)|=(x-u)\,(|p|_a-m|p|_b).\]
The hypothesis gives a proper nonempty prefix \(p\) with \(|p|_a=m|p|_b\) and \(|p|_b\in\{1,2\}\). Then \(g(p),h(p)\) are prefixes of the same word of equal length, so \(p\in\mathrm{Eq}(g,h)\); writing \(w=pq\) and cancelling gives \(q\in\mathrm{Eq}(g,h)\).

If \(\mathrm{Eq}(g,h)=z^*\), then \(p,w\in z^*\) and \(|z|_b\) divides \(|p|_b\in\{1,2\}\) and \(3\); so \(z=a^rba^s\) and \(w=z^3\in A\).

Otherwise \(\mathrm{Eq}(g,h)\) has at least two generators, and by Holub's structure theorem \cite{uniq} (using \cite{gen2}) it equals \(\{a^nb,ba^n\}^*\) with \(n\ge1\), up to exchanging the letters. In the unexchanged case \(w\) is a product of three generators and \(3n=S\), so \(n=m\). With \(L=a^mb\), \(R=ba^m\) the eight products give
\[
LLL,RRR\in A;\quad LLR=(m,m,0,m),\ LRR=(m,0,m,m)\in C;\quad LRL,RLL,RLR,RRL\in D .
\]
In the exchanged case \(\{b^na,ab^n\}^*\), \(w\) is a product of \(c\) generators with \(cn=3\). The case \(c=1\) is impossible because \(w\) has the proper nonempty prefix \(p\) in \(\mathrm{Eq}(g,h)\). So \(n=1\) and \(w\) is a product of three words from \(\{ab,ba\}\), which is the list above with \(m=1\).
\end{proof}

\begin{corollary}\label{outer}
For \(i,\ell\ge1\), the tuple \((i,0,0,\ell)\) forces periodicity under the length condition. Together with family \(B\), the slice \(j=k=0\) is completely classified.
\end{corollary}
\begin{proof}
For \(i+\ell\ge3\) this is \cite[Theorem 1]{hh} (the word \(a^ib^3a^\ell\)), with Lemma~\ref{perred} covering periodic morphisms. The remaining case \((1,0,0,1)\) is Theorem~\ref{t1001}. If \(i=0\) or \(\ell=0\), the tuple lies in \(B\).
\end{proof}

\begin{lemma}[Padding preserves forcing]\label{padding}
If \((i,j,k,\ell)\) forces periodicity under the length condition, so does \((i+r,\,j+r+s,\,k+r+s,\,\ell+s)\) for all \(r,s\ge0\).
\end{lemma}
\begin{proof}
A solution \(X,Y,U,V\) of the padded equation gives the solution \(X,\ Y'=X^rYX^s,\ U,\ V'=U^rVU^s\) of the original one. The condition \(|V|\ge|Y|\) forces \(x\ge u\), so \(|V'|\ge|V|\ge|X|,|U|\), and \(|V'|-|Y'|=\tfrac{S}{3}(x-u)\ge0\), where \(S\) is the original exponent sum. If \(X=U\) and \(Y'=V'\), cancelling powers of \(X\) gives \(Y=V\). Otherwise \(X,Y',U,V'\) are powers of one primitive word, hence so are \(Y\) and \(V\).
\end{proof}
Among the \(625\) tuples with exponents at most \(4\), Theorem~\ref{balanced} applies to \(178\). The \(30\) of them in \(A\cup C\cup D\) are exactly the nonforcing ones found by the search, so the remaining \(148\) are now proved forcing for all lengths. The open cases of the conjecture are the tuples outside the families with no balanced prefix, for example \(2011\) and \(1001\), which needed the arguments of Sections~2 and~5. By Lemma~\ref{perred}, any remaining counterexample has \(XY\ne YX\) and \(UV\ne VU\).

\section{Infinite slices proved by periodicity arguments}
Throughout, \(x,y,u,v\) are the lengths of \(X,Y,U,V\), and \(W\) is the common word. We use the Fine--Wilf theorem: a word with periods \(p,q\) and length at least \(p+q-\gcd(p,q)\) has period \(\gcd(p,q)\). We also use two standard facts about a primitive word \(R\) with \(r=|R|\). (P1) If a factor of \(R^\omega\) of length at least \(r+c\) has period \(c\), then \(r\mid c\): by Fine--Wilf it has period \(\gcd(r,c)\), and a proper divisor would make a conjugate of \(R\) imprimitive. (P2) A factor of \(R^\omega\) whose length is a multiple of \(r\) is a power of a conjugate of \(R\), so any word \(Z\) with \(Z^m\) such a factor and \(Z\ne\eps\) has \(r\mid |Z|\). In each proof, if \(x=u\) the length equation gives \(y=v\) and the substitutions coincide, so we assume \(x\ne u\), and by exchanging the substitutions, \(x>u\).

\begin{theorem}\label{ijjl}
Let \(0\le i,\ell\le j\), \(j\ge2\) and \(i+\ell\ne j\). Then \((i,j,j,\ell)\) forces periodicity, even without the length condition.
\end{theorem}
\begin{proof}
Let \(\delta=x-u\). Then \((i+2j+\ell)\delta=3(v-y)\). Put \(P=jx+y\), \(Q=ju+v\). Since \(i,\ell\le j\), \(W\) is a factor of \((X^jY)^\omega\) and of \((U^jV)^\omega\), so it has periods \(P\) and \(Q\), and \(P-Q=\tfrac{j-i-\ell}{3}\delta\ne0\). As \(|W|\ge2P,2Q\), Fine--Wilf gives the period \(\gcd(P,Q)\). Let \(R\) be the primitive root of the prefix of that length and \(p=|R|\). Then \(p\mid P,Q\) and \(p\le|P-Q|\le j\delta/3\le jx/3\).

The factor \(X^j\) has periods \(x\) and \(p\), and \(jx\ge x+p\) because \((j-1)x\ge jx/2\ge p\). So it has period \(\gcd(x,p)\), and primitivity of \(R\) gives \(p\mid x\), hence \(p\mid y\). Every left block boundary is a multiple of \(p\), so \(X,Y\in R^*\) and \(p\mid|W|\). Now \(|W|=3Q+(i+\ell-j)u\) gives \(p\mid mu\) with \(m=|i+\ell-j|\in[1,j]\). The factor \(U^m\) of \(U^j\) has length divisible by \(p\), so \(p\mid u\) by (P2), and \(p\mid v\). All right block boundaries are multiples of \(p\), so \(U,V\in R^*\).
\end{proof}

\begin{theorem}\label{0jk0}
For \(j,k\ge1\), the tuple \((0,j,k,0)\) forces periodicity, even without the length condition, unless \(j=2k\), \(k=2j\) or \(j=k=1\).
\end{theorem}
\begin{proof}
By reversal assume \(j\le k\). Let \(\delta=x-u>0\), so \(3(v-y)=(j+k)\delta\). Put \(A=YX^j\), \(B=VU^j\), \(P=|A|\), \(Q=|B|\). Then \(Q-P=\tfrac{k-2j}{3}\delta\), and \(A^2\), \(B^2\) are both prefixes of \(W\). The case \(j=k\ge2\) is Theorem~\ref{ijjl}.

\emph{Case \(j<k<2j\).} Then \(j\ge2\). Put \(h=P-Q=\tfrac{2j-k}{3}\delta\), so \(0<h\le(j-1)x/3\). Write \(A=BC\) with \(|C|=h\). From \(BB\preceq BCBC\) we get \(B\preceq CBC\), so \(B\) has period \(h\) and \(B\preceq R^\omega\), where \(R\) is the primitive root of \(C\) and \(r=|R|\mid h\). The word \(B\) contains a prefix of \(X^j\) of length \(jx-h\ge x+r\). By Fine--Wilf and primitivity, \(r\mid x\), and that prefix of \(X^j\) is a power of the rotation of \(R\) by \(y\bmod r\). The suffix of \(X^j\) of length \(h\) is \(C=R^{h/r}\), so this rotation is \(R\) itself, and \(r\mid y\). Hence \(X,Y\in R^*\) and \(r\mid P,Q\). From \(|W|=3Q+(k-2j)u\) we get \(r\mid(2j-k)u\) with \(1\le 2j-k<j\), so (P2) gives \(r\mid u\) and \(r\mid v\). All four words lie in \(R^*\).

\emph{Case \(k>2j\).} Put \(h=Q-P>0\) and \(F=YX^k\), so \(W=AFY\). Since \(B^2\) is a prefix of \(W\), \(W[s:t]=W[Q+s:Q+t]\) for \(0\le s\le t\le Q\). With \(s=y\), \(t=Q\) the left side is \(X^jF[0:h]\). The right side starts \(h\) letters into \(X^k\) and has length \(jx+h\), and \(jx+2h\le kx\), so
\[X^jF[0:h]=X^k[h:jx+2h].\]
Comparing the first \(x\) letters (\(j\ge1\)), the rotation of \(X\) by \(h\) equals \(X\), so \(r\mid h\), where \(R\) is the primitive root of \(X\). Hence the right side is \(X^jH\) with \(H=R^{h/r}\), and \(F[0:h]=H\). With \(s=0\), \(t=y\) we get \(Y=F[h:h+y]\), so \(F[0:h+y]=HY\). But \(F=YX^k\) begins with \(YH\). Thus \(HY=YH\), and since \(H\) is a nonempty power of \(R\), \(Y\in R^*\). Then \(r\mid P,Q\), and \(|W|=3Q+(k-2j)u\) with \(1\le k-2j\le k\) gives \(r\mid u\) and \(r\mid v\) by (P2).
\end{proof}

\begin{theorem}\label{0020}
The tuple \((0,0,2,0)\) forces periodicity.
\end{theorem}
\begin{proof}
The equation is \(Y^2X^2Y=V^2U^2V\). Lengths give \(2(x-u)=3(v-y)\), and \(v\ge y\) forces \(x=u+3s\), \(v=y+2s\) with \(s\ge0\). If \(s=0\) the substitutions coincide, so let \(s>0\). The condition \(v\ge x\) gives \(y\ge u+s\). Write \(V=YA=DY\) with \(|A|=|D|=2s\).

If \(y\ge2s\), then \(V\preceq Y^2\), so \(A\) and \(D\) both equal the prefix of \(Y\) of length \(2s\). Hence \(YA=AY\), and \(Y\), \(V\) commute.

If \(y<2s\), put \(e=2s-y\). Then \(0<e\le s\), \(u\le s-e\), and \(V=YZY\) with \(|Z|=e\). Comparing \(Y^2\) with \(V^2\) gives \(ZY=YT\) with \(|T|=e\), so \(Y\) has period \(e\) and \(V=Y^2T\). Cancelling \(Y^2\) and the final \(Y\) gives \(X^2=TY^2TU^2YZ\). Reading the two copies of \(X\),
\[TY\,Y[0:s+u]=LU^2YZ,\]
where \(L\) is the suffix of \(YT\) of length \(s-u\). The marked occurrences of \(Y\) start at \(e\) and at \(s+u\). Put \(d=s+u-e\ge0\). If \(d=0\), then \(u=0\), \(e=s=y\), and \(ZY=YT\) gives \(Z=Y=T\). Otherwise \(Y\) has periods \(d\) and \(e\), and \(y\ge d+e\), so it has period \(g=\gcd(d,e)\). The overlap of \(Y[0:s+u]\) on the left with the occurrence of \(Y\) on the right gives \(Y[q]=Y[s-u+q]\) for \(0\le q<d\). Since \(g\mid e\) and \(g\mid s+u\), we have \(s-u\equiv 2s\equiv y\pmod g\). So the prefix \(Y[0:g]\) is fixed by rotation by \(y\bmod g\). Its primitive root \(\rho\) therefore satisfies \(|\rho|\mid y\), and \(Y\) is a power of \(\rho\). The words \(Z\) and \(T\) are a prefix and a suffix of \(Y\) of length \(e\), a multiple of \(|\rho|\), so they are powers of \(\rho\) too, and \(V=YZY\) commutes with \(Y\).

In both cases \(Y=R^a\), \(V=R^b\) with \(R\) primitive and \(b>a\). Cancelling gives \(X^2=H^2U^2H\) with \(H=R^{b-a}\ne\eps\). If \(HU\ne UH\), then \(\{H,U\}\) is a binary code, and every primitive word over it whose image is imprimitive is conjugate to a word in \(a^*b\cup ab^*\) (see \cite{hh}, Section~2). The word \(a^2b^2a\) is primitive and not of that form, so \(H^2U^2H\) would be primitive, contradicting \(X\ne\eps\). So \(U\in R^*\), and then \(X\in R^*\).
\end{proof}

\begin{corollary}\label{0jk0full}
For all \(j,k\ge0\), the tuple \((0,j,k,0)\) is nonforcing exactly when \((j,k)\in\{(0,0),(0,1),(1,0),(1,1)\}\), or \(j,k>0\) and \(j=2k\) or \(k=2j\).
\end{corollary}
\begin{proof}
For \(j,k>0\) this is Theorem~\ref{0jk0} with families \(D\) and \(0110\). By reversal let \(j=0\). The tuple \((0,0,k,0)\) is the word \(b^2a^kb\), so for \(k\ge3\) it forces periodicity by \cite[Theorem~1]{hh}, which has no length condition and covers periodic morphisms through its Lemma~9. The case \(k=2\) is Theorem~\ref{0020}, and \(k=0,1\) are \(0000\in A\) and the sporadic \(0010\).
\end{proof}

\begin{theorem}\label{ijj0}
If \(i>j\ge1\), then \((i,j,j,0)\) forces periodicity, even without the length condition. By reversal, so does \((0,j,j,\ell)\) for \(\ell>j\).
\end{theorem}
\begin{proof}
Put \(n=i-j\), \(A=X^jY\), \(B=U^jV\), \(P=|A|\), \(Q=|B|\). The equation is \(X^nA^3=U^nB^3\), and \(Q-P=n\delta/3>0\) with \(\delta=x-u\). The factor \(X^i\) overlaps the suffix \(B^3\) in a word of length \(ix-nu=jx+3(Q-P)\).

If \(Q\le2P\), the suffix \(A^3\) of \(B^3\) has length \(3P\ge P+Q\), so by Fine--Wilf the primitive roots of \(A\) and \(B\) are conjugate, of a common length \(r\) with \(r\mid P,Q\) and \(r\le Q-P\). The overlap has length at least \(x+r\), so the primitive root of \(X\) is conjugate to them as well.

If \(Q>2P\), the overlap has length greater than \(x+Q\), so the primitive roots of \(X\) and \(B\) are conjugate, of length \(r\le x\le P\). Fine--Wilf on the suffix \(A^3\) brings in the root of \(A\).

In both cases \(X\) and \(A=X^jY\) start at the same phase, so they are powers of one primitive word \(R\), and \(Y\in R^*\). Then \(W\in R^*\) and \(r\mid Q\), so \(r\mid |W|-3Q=nu\). The prefix \(U^n\) is a power of \(R\), so \(U\in R^*\), and then \(V\in R^*\).
\end{proof}

\begin{corollary}\label{ijjlfull}
The tuple \((i,j,j,\ell)\) is nonforcing exactly when it is in family \(A\) (\(i+\ell=j\)), in family \(B\) (\(j=0\) and \(i\ell=0\)), or equals \(0110\) or \(1111\).
\end{corollary}
\begin{proof}
For \(j\ge2\) and \(i,\ell\le j\), use Theorem~\ref{ijjl}. If one outer exponent is \(0\) and the other exceeds \(j\ge1\), use Theorem~\ref{ijj0}. If \(i,\ell>0\) and \(\max(i,\ell)>j\), choose \(r+s=j\) with \(r<i\), \(s<\ell\). With \(Y'=X^rYX^s\) and \(V'=U^rVU^s\), the tuple is a padding of \((i-r,0,0,\ell-s)\), which forces periodicity by Corollary~\ref{outer}, so Lemma~\ref{padding} applies. The cases \(j\le1\) with \(i,\ell\le j\), and \(j=0\), are the listed exceptions or Corollary~\ref{outer}.
\end{proof}

\begin{lemma}\label{single}
Suppose \(P^2XP=Q^2UQ\), where \(X\) is a prefix or a suffix of \(P\) and \(U\) is a prefix or a suffix of \(Q\). Then \(X=U\) and \(P=Q\), or all four words commute.
\end{lemma}
\begin{proof}
Lengths give \(x-u=3(|Q|-|P|)\). By symmetry let \(h=|Q|-|P|>0\). Then \(|P|\ge x\ge3h\), and \(P\) is both a prefix and a suffix of \(Q\), so \(Q=PA=DP\) with \(|A|=|D|=h\). Because \(|Q|\le2|P|\), both \(A\) and \(D\) are the prefix of \(P\) of length \(h\). So \(PA=AP\), and \(P,Q,A\in R^*\) for a primitive \(R\). Cancelling gives \(X=A^2UA\). Since \(X\) is a prefix or suffix of \(P\in R^*\) and begins and ends with the nonempty power \(A\), we get \(|R|\mid x\), so \(X\in R^*\), and then \(U\in R^*\).
\end{proof}

\begin{theorem}\label{region}
If \(\min(j,k)\ge i+\ell\), the tuple \((i,j,k,\ell)\) forces periodicity unless it lies in \(A\cup C\cup D\) or is sporadic.
\end{theorem}
\begin{proof}
Let \(c=i+\ell\). With \(Y'=X^iYX^\ell\), the tuple is a padding of \((0,j-c,k-c,0)\). If the reduced tuple forces periodicity, Lemma~\ref{padding} finishes. By Corollary~\ref{0jk0full}, the remaining reduced tuples are: \((0,0)\), which gives family \(A\); \((1,1)\), covered by Corollary~\ref{ijjlfull}; \(j-c=2(k-c)>0\) or the reverse, where \(S=3(c+\min)\) and Theorem~\ref{balanced} applies; and \((0,1)\) or \((1,0)\). In the last case the equation is \(P^2XP=Q^2UQ\) or its reversal, with \(P=X^iYX^\ell\), and Lemma~\ref{single} applies when \(c>0\). If \(c=0\) the tuple is the sporadic \(0010\) or \(0100\).
\end{proof}

\begin{remark}
If \(\min(j,k)<i+\ell\) and \(j\le k\), choose \(r+s=j\) with \(r\le i\), \(s\le\ell\). Padding reduces \((i,j,k,\ell)\) to \((i-r,0,k-j,\ell-s)\). So the full conjecture reduces to classifying the family \((i,0,k,\ell)\), that is \(X^iY^2X^kYX^\ell=U^iV^2U^kVU^\ell\), together with the padded descendants of its exceptions, since padding only transfers forcing. All results of this section agree with \texttt{map4.txt} and with the conjectured list for every exponent up to \(29\).
\end{remark}

\section{Verification and priority}
Lean 4.33.0 checked \texttt{Holub.forcing} and \texttt{Holub.words2025} over arbitrary letter types. The first has no length hypotheses. The second spells out the seven blocks and all three inequalities. Both use only \texttt{propext} and \texttt{Quot.sound} according to \texttt{\#print axioms}. Theorem~\ref{t1001} is checked as \texttt{Holub1001.case1001} in \texttt{Case1001.lean}, also using only \texttt{propext} and \texttt{Quot.sound}. Both compiled modules were replayed with \texttt{leanchecker}. Corollary~\ref{gcd}, the exponent families of Section~4, Theorem~\ref{nonforcing} and Sections~7--8 have written proofs only; the nonforcing witnesses are also checked by \texttt{families.py}. Independently of Lean, the disjoint-set scan covered all \(1{,}703{,}400\) admissible \(t>0\) tuples with \(|V|\le320\), including \(1{,}419{,}156\) primitive scales. Each primitive graph was connected. Tests compare the general graph solver against literal binary substitution and check transport, boundary signs, and examples.

The conference collection dated 1 July 2025 explicitly posed the target on p.~5. Searches through 24 November 2026 checked Holub's publications page, primary papers, indexed proceedings and preprints, using the literal patterns and the terms \emph{binary equality words}, \emph{three occurrences} and \emph{periodicity forcing}. Hadravov\'a and Holub's Theorem 1 concerns a single internal run of \(b\)'s with both outer exponents positive. Holub and S\'ykora's Theorem 30 classifies equality words with two \(b\)'s for two nonperiodic morphisms. Neither result covers the target. No earlier resolution of this pattern was located. Bibliographic priority and external expert review remain unverified.

\paragraph{Contribution disclosure.}
The author selected the problem from the WORDS 2025 collection, directed the investigation, and audited and independently re-verified the results, including compiling the Lean proof with Lean~4.33.0 and \texttt{leanchecker}, rerunning the exhaustive search, and checking small cases by direct enumeration. The AI system Astra~6.2 developed the written proof, implemented the searches, wrote the Lean formalization and drafted the manuscript under this direction. The reductions and periodicity arguments of Sections~7--8 came from an AI-assisted review and were checked line by line and against the search data.

\begin{thebibliography}{99}
\bibitem{words}\v S. Holub, Periodicity forcing, WORDS 2025 Open Problems, Nancy, 1 July 2025, slides 3--5. \href{https://words2025.loria.fr/files/2025/07/webpage-open-problems.pdf}{Conference collection}.
\bibitem{hh}J. Hadravov\'a and \v S. Holub, \emph{Equation \(x^iy^jx^k=u^iv^ju^k\) in words}, 15 January 2015, Theorem 1. \href{https://arxiv.org/abs/1501.03602}{arXiv:1501.03602}.
\bibitem{hs}\v S. Holub and J. S\'ykora, \emph{Binary equality words with two b's}, Comment. Math. Univ. Carolin. 59(2) (2018), 153--172. \href{https://cmuc.karlin.mff.cuni.cz/pdf/cmuc1802/holsyk.pdf}{Article}.
\bibitem{periodic}\v S. Holub, \emph{Binary equality words for periodic morphisms}, author version, Theorem 1; published as \emph{Binary equality languages for periodic morphisms}, RIMS K\^oky\^uroku 1366 (2004). \href{https://www.karlin.mff.cuni.cz/~holub/soubory/rimsrevised.pdf}{Author manuscript}.
\bibitem{gen2}\v S. Holub, \emph{Binary equality sets are generated by two words}, J. Algebra 259 (2003), 1--42. \href{https://doi.org/10.1016/S0021-8693(02)00534-3}{Article}.
\bibitem{uniq}\v S. Holub, \emph{A unique structure of two-generated binary equality sets}, Developments in Language Theory (DLT 2002), LNCS 2450 (2003), 245--257. \href{https://www.karlin.mff.cuni.cz/~holub/soubory/uniq.pdf}{Author version}.
\bibitem{mak}G. S. Makanin, \emph{The problem of solvability of equations in a free semigroup}, Math. USSR-Sb. 32(2) (1977), 129--198. \href{https://doi.org/10.1070/SM1977v032n02ABEH002376}{Article}.
\end{thebibliography}
\end{document}
